\documentclass[11pt,a4paper]{article}

% -------------------------------------------------
% Packages
% -------------------------------------------------
\usepackage[margin=2.2cm]{geometry}
\usepackage{kotex}
\usepackage{amsmath,amssymb}
\usepackage{booktabs}
\usepackage{array}
\usepackage{xcolor}
\usepackage{listings}
\usepackage[most]{tcolorbox}
\usepackage{enumitem}
\usepackage{hyperref}
\usepackage{fancyhdr}
\usepackage{titlesec}

% -------------------------------------------------
% Colors
% -------------------------------------------------
\definecolor{codebg}{RGB}{247,247,247}
\definecolor{codeframe}{RGB}{210,210,210}
\definecolor{keywordcolor}{RGB}{0,70,160}
\definecolor{commentcolor}{RGB}{0,120,80}
\definecolor{stringcolor}{RGB}{170,50,50}
\definecolor{titleblue}{RGB}{35,75,130}

% -------------------------------------------------
% Hyperlinks
% -------------------------------------------------
\hypersetup{
    colorlinks=true,
    linkcolor=titleblue,
    urlcolor=titleblue
}

% -------------------------------------------------
% Header / Footer
% -------------------------------------------------
\pagestyle{fancy}
\setlength{\headheight}{14pt}
\fancyhf{}
\lhead{C++ Programming}
\rhead{Day 2}
\cfoot{\thepage}

% -------------------------------------------------
% Section style
% -------------------------------------------------
\titleformat{\section}
  {\Large\bfseries\color{titleblue}}
  {\thesection.}{0.6em}{}

\titleformat{\subsection}
  {\large\bfseries}
  {\thesubsection.}{0.5em}{}

% -------------------------------------------------
% C++ code style
% -------------------------------------------------
\lstdefinestyle{cppstyle}{
    language=C++,
    backgroundcolor=\color{codebg},
    frame=single,
    rulecolor=\color{codeframe},
    basicstyle=\ttfamily\small,
    keywordstyle=\color{keywordcolor}\bfseries,
    commentstyle=\color{commentcolor},
    stringstyle=\color{stringcolor},
    numbers=left,
    numberstyle=\tiny\color{gray},
    numbersep=8pt,
    tabsize=4,
    showstringspaces=false,
    breaklines=true,
    columns=fullflexible
}
\lstset{style=cppstyle}

% -------------------------------------------------
% Boxes
% -------------------------------------------------
\newtcolorbox{learningbox}{
    colback=blue!4,
    colframe=titleblue,
    title=\textbf{학습 목표},
    fonttitle=\bfseries
}

\newtcolorbox{importantbox}{
    colback=yellow!8,
    colframe=orange!70!black,
    title=\textbf{핵심 포인트},
    fonttitle=\bfseries
}

\newtcolorbox{checkpointbox}{
    colback=green!5,
    colframe=green!45!black,
    title=\textbf{체크 질문},
    fonttitle=\bfseries
}

\newtcolorbox{practicebox}{
    colback=red!3,
    colframe=red!55!black,
    title=\textbf{실습},
    fonttitle=\bfseries
}

% -------------------------------------------------
% Document
% -------------------------------------------------
\begin{document}

\begin{titlepage}
    \centering
    \vspace*{3cm}

    {\Huge\bfseries C++ Programming\par}
    \vspace{0.8cm}
    {\LARGE Day 2: 조건 분기와 반복 제어\par}

    \vspace{2cm}

    {\large if--else, switch, do--while, break, continue, 난수\par}

    \vfill
\end{titlepage}

\tableofcontents
\newpage

% =================================================
\section{학습 목표}
% =================================================

\begin{learningbox}
이 장을 마치면 다음 내용을 설명하고 직접 코드로 작성할 수 있다.
\begin{itemize}[leftmargin=1.5em]
    \item \texttt{if--else}와 \texttt{else if}를 이용한 다중 조건 분기
    \item 중첩 조건문을 이용한 단계적 판단
    \item \texttt{switch}문과 \texttt{break}의 역할
    \item \texttt{while}문과 \texttt{do--while}문의 차이
    \item \texttt{break}, \texttt{continue}를 이용한 반복 흐름 제어
    \item \texttt{rand()}, \texttt{srand()}, \texttt{time()}을 이용한 난수 생성
    \item 조건문, 반복문, 난수를 결합한 숫자 맞히기 게임 구현
\end{itemize}
\end{learningbox}

Day 1에서는 프로그램의 기본 구조와 기본적인 조건문 및 반복문을 학습했다.
Day 2에서는 이미 배운 문법을 더 정교하게 사용하여 프로그램의 실행 흐름을 설계한다.
중요한 것은 문법을 외우는 것이 아니라, 어떤 조건에서 어느 코드가 실행되고 반복이 언제 계속되거나 종료되는지를 정확히 이해하는 것이다.

% =================================================
\section{두 가지 경우 선택: \texttt{if--else}}
% =================================================

\texttt{if}문은 조건이 참일 때만 코드를 실행한다. 그러나 실제 프로그램에서는 조건이 참일 때와 거짓일 때 서로 다른 작업을 수행해야 하는 경우가 많다. 이때 \texttt{if--else}를 사용한다.

\begin{lstlisting}
if (condition)
{
    // condition이 true일 때 실행
}
else
{
    // condition이 false일 때 실행
}
\end{lstlisting}

\subsection{예제: 성인 여부 판별}

\begin{lstlisting}
#include <iostream>

using namespace std;

int main()
{
    int age;

    cout << "Age: ";
    cin >> age;

    if (age >= 20)
    {
        cout << "Adult" << endl;
    }
    else
    {
        cout << "Minor" << endl;
    }

    return 0;
}
\end{lstlisting}

입력이 23이면 \texttt{age >= 20}이 참이므로 \texttt{Adult}가 출력된다. 입력이 17이면 조건이 거짓이므로 \texttt{Minor}가 출력된다.

\begin{importantbox}
\texttt{if--else}에서는 두 블록 중 정확히 하나만 실행된다. 서로 동시에 성립할 수 없는 두 경우를 처리할 때 적합하다.
\end{importantbox}

\subsection{비교 연산자와 대입 연산자}

다음 두 표현은 의미가 완전히 다르다.

\begin{lstlisting}
number = 10;   // 대입
number == 10;  // 비교
\end{lstlisting}

\texttt{=}는 오른쪽 값을 왼쪽 변수에 저장하는 대입 연산자이고, \texttt{==}는 두 값이 같은지 검사하는 비교 연산자이다.

\begin{lstlisting}
if (number == 10)
{
    cout << "The number is 10." << endl;
}
\end{lstlisting}

조건식에서 실수로 \texttt{number = 10}을 쓰면 10이 변수에 대입되고, 0이 아닌 값이 참으로 평가되어 의도와 다르게 작동할 수 있다.

\begin{practicebox}
\textbf{파일: \texttt{01\_if\_else.cpp}}\\
정수 하나를 입력받아 0 이상이면 \texttt{Non-negative}, 0보다 작으면 \texttt{Negative}를 출력하라.
\end{practicebox}

% =================================================
\section{여러 조건 처리: \texttt{else if}}
% =================================================

두 가지보다 많은 경우를 처리하려면 \texttt{else if}를 사용한다.

\begin{lstlisting}
if (condition1)
{
    // condition1이 true
}
else if (condition2)
{
    // condition1은 false이고 condition2는 true
}
else
{
    // 앞의 모든 조건이 false
}
\end{lstlisting}

\subsection{예제: 학점 판별}

\begin{lstlisting}
#include <iostream>

using namespace std;

int main()
{
    int score;

    cout << "Score: ";
    cin >> score;

    if (score >= 90)
    {
        cout << "Grade A" << endl;
    }
    else if (score >= 80)
    {
        cout << "Grade B" << endl;
    }
    else if (score >= 70)
    {
        cout << "Grade C" << endl;
    }
    else if (score >= 60)
    {
        cout << "Grade D" << endl;
    }
    else
    {
        cout << "Grade F" << endl;
    }

    return 0;
}
\end{lstlisting}

\subsection{조건은 위에서 아래로 검사된다}

\texttt{else if} 구조에서는 조건을 위에서부터 차례로 검사한다. 처음으로 참이 된 블록 하나만 실행되고 나머지는 검사하지 않는다.

따라서 다음처럼 낮은 기준을 먼저 쓰면 잘못된 결과가 나온다.

\begin{lstlisting}
if (score >= 60)
{
    cout << "Grade D" << endl;
}
else if (score >= 90)
{
    cout << "Grade A" << endl;
}
\end{lstlisting}

점수가 95여도 첫 번째 조건이 이미 참이므로 \texttt{Grade D}가 출력된다. 범위가 겹치는 조건은 일반적으로 더 제한적인 조건부터 작성해야 한다.

\begin{importantbox}
\texttt{else if}는 위에서부터 검사하고, 처음 참인 블록 하나만 실행한다. 따라서 조건의 순서가 프로그램의 결과를 바꿀 수 있다.
\end{importantbox}

\begin{practicebox}
\textbf{파일: \texttt{02\_else\_if.cpp}}\\
정수 하나를 입력받아 양수이면 \texttt{Positive}, 0이면 \texttt{Zero}, 음수이면 \texttt{Negative}를 출력하라.
\end{practicebox}

% =================================================
\section{조건문 안의 조건문: Nested \texttt{if}}
% =================================================

조건문 내부에 또 다른 조건문을 작성하는 것을 중첩 조건문이라고 한다.

\begin{lstlisting}
if (condition1)
{
    if (condition2)
    {
        // 두 조건이 모두 참일 때 실행
    }
}
\end{lstlisting}

\subsection{예제: 단계적 로그인 검사}

\begin{lstlisting}
#include <iostream>
#include <string>

using namespace std;

int main()
{
    string password;
    bool isAdmin = true;

    cout << "Password: ";
    cin >> password;

    if (password == "cpp123")
    {
        cout << "Login successful." << endl;

        if (isAdmin)
        {
            cout << "Administrator mode." << endl;
        }
    }
    else
    {
        cout << "Login failed." << endl;
    }

    return 0;
}
\end{lstlisting}

안쪽의 \texttt{if (isAdmin)}은 바깥쪽 비밀번호 검사가 성공했을 때만 평가된다.

두 조건을 항상 동시에 검사해도 된다면 논리 연산자 \texttt{\&\&}를 사용할 수 있다.

\begin{lstlisting}
if (password == "cpp123" && isAdmin)
{
    cout << "Administrator mode." << endl;
}
\end{lstlisting}

단계별로 다른 메시지를 출력해야 한다면 중첩 조건문이 더 자연스럽다. 다만 중첩이 지나치게 깊으면 코드 흐름을 이해하기 어려워지므로 주의한다.

\begin{practicebox}
\textbf{파일: \texttt{03\_nested\_if.cpp}}\\
정수 하나를 입력받아 먼저 양수인지 검사하라. 양수라면 다시 짝수인지 홀수인지 검사하고, 0 또는 음수라면 \texttt{Not a positive number}를 출력하라.
\end{practicebox}

% =================================================
\section{하나의 값으로 여러 경우 선택: \texttt{switch}}
% =================================================

하나의 변수나 표현식이 특정 값과 일치하는지 여러 번 비교할 때 \texttt{switch}문을 사용할 수 있다.

\begin{lstlisting}
switch (expression)
{
    case value1:
        // expression == value1
        break;

    case value2:
        // expression == value2
        break;

    default:
        // 일치하는 case가 없을 때
        break;
}
\end{lstlisting}

\subsection{예제: 메뉴 선택}

\begin{lstlisting}
#include <iostream>

using namespace std;

int main()
{
    int menu;

    cout << "1. Start" << endl;
    cout << "2. Settings" << endl;
    cout << "3. Exit" << endl;
    cout << "Select: ";
    cin >> menu;

    switch (menu)
    {
        case 1:
            cout << "Game started." << endl;
            break;

        case 2:
            cout << "Settings opened." << endl;
            break;

        case 3:
            cout << "Program ended." << endl;
            break;

        default:
            cout << "Invalid menu." << endl;
            break;
    }

    return 0;
}
\end{lstlisting}

\subsection{\texttt{break}와 fall-through}

\texttt{break}는 현재 \texttt{switch}문을 즉시 종료한다. \texttt{break}가 없으면 일치한 \texttt{case} 아래의 코드가 다음 \texttt{case}까지 연속해서 실행될 수 있다.

\begin{lstlisting}
int number = 1;

switch (number)
{
    case 1:
        cout << "One" << endl;
    case 2:
        cout << "Two" << endl;
    case 3:
        cout << "Three" << endl;
}
\end{lstlisting}

출력 결과:

\begin{verbatim}
One
Two
Three
\end{verbatim}

이 현상을 fall-through라고 한다. 의도적으로 사용할 수도 있지만, 초반에는 각 \texttt{case}의 마지막에 \texttt{break}를 작성하는 것이 안전하다.

\subsection{\texttt{if--else}와 \texttt{switch} 비교}

\begin{center}
\begin{tabular}{>{\raggedright\arraybackslash}p{0.45\textwidth} >{\raggedright\arraybackslash}p{0.3\textwidth}}
\toprule
상황 & 적합한 문법 \\
\midrule
값의 범위를 검사 & \texttt{if--else} \\
두 조건을 동시에 검사 & \texttt{if--else} \\
하나의 값과 여러 상수를 비교 & \texttt{switch} \\
메뉴 번호 선택 & \texttt{switch} \\
복잡한 논리 조건 & \texttt{if--else} \\
\bottomrule
\end{tabular}
\end{center}

\begin{practicebox}
\textbf{파일: \texttt{04\_switch.cpp}}\\
1부터 7 사이의 숫자를 입력받아 해당 요일의 영어 이름을 출력하라. 그 외의 값은 \texttt{Invalid day}로 처리한다.
\end{practicebox}

% =================================================
\section{최소 한 번 실행: \texttt{do--while}}
% =================================================

\texttt{while}문은 조건을 먼저 검사한다. 따라서 처음부터 조건이 거짓이면 본문이 한 번도 실행되지 않는다.

\begin{lstlisting}
int number = 10;

while (number < 5)
{
    cout << number << endl;
}
\end{lstlisting}

반면 \texttt{do--while}문은 본문을 먼저 실행한 뒤 조건을 검사한다.

\begin{lstlisting}
do
{
    // 반복할 코드
}
while (condition);
\end{lstlisting}

\begin{lstlisting}
#include <iostream>

using namespace std;

int main()
{
    int number = 10;

    do
    {
        cout << number << endl;
    }
    while (number < 5);

    return 0;
}
\end{lstlisting}

조건은 거짓이지만 본문이 먼저 실행되므로 10이 한 번 출력된다.

\subsection{입력 검증}

사용자의 입력을 먼저 받아야 조건을 검사할 수 있으므로, 입력 검증에는 \texttt{do--while}문이 자연스럽다.

\begin{lstlisting}
int number;

do
{
    cout << "Enter a number from 1 to 10: ";
    cin >> number;
}
while (number < 1 || number > 10);

cout << "Accepted: " << number << endl;
\end{lstlisting}

\begin{importantbox}
\texttt{do--while}문의 마지막 \texttt{while (condition)} 뒤에는 반드시 세미콜론이 필요하다.
\end{importantbox}

\begin{practicebox}
\textbf{파일: \texttt{05\_do\_while.cpp}}\\
사용자가 0을 입력할 때까지 정수를 계속 입력받아 각 값을 출력하라. 0을 입력하면 반복을 종료한다.
\end{practicebox}

% =================================================
\section{반복 종료: \texttt{break}}
% =================================================

\texttt{break}는 현재 실행 중인 가장 가까운 반복문이나 \texttt{switch}문을 즉시 종료한다.

\begin{lstlisting}
for (int number = 1; number <= 10; number++)
{
    if (number == 5)
    {
        break;
    }

    cout << number << endl;
}
\end{lstlisting}

출력 결과:

\begin{verbatim}
1
2
3
4
\end{verbatim}

\subsection{무한 반복과 \texttt{break}}

반복 횟수를 미리 정하기 어려울 때는 무한 반복을 만든 뒤 특정 조건에서 종료할 수 있다.

\begin{lstlisting}
while (true)
{
    int number;

    cout << "Enter 0 to quit: ";
    cin >> number;

    if (number == 0)
    {
        break;
    }

    cout << "You entered " << number << endl;
}
\end{lstlisting}

중첩 반복문 안의 \texttt{break}는 가장 가까운 반복문 하나만 종료한다.

\begin{practicebox}
\textbf{파일: \texttt{06\_break.cpp}}\\
1부터 100까지 차례로 검사하다가 처음으로 7의 배수이면서 20보다 큰 수를 찾으면 출력하고 반복을 종료하라.
\end{practicebox}

% =================================================
\section{현재 반복 건너뛰기: \texttt{continue}}
% =================================================

\texttt{continue}는 반복문 전체를 종료하지 않는다. 현재 반복에서 남은 코드를 건너뛰고 다음 반복으로 이동한다.

\begin{lstlisting}
for (int number = 1; number <= 10; number++)
{
    if (number % 2 == 0)
    {
        continue;
    }

    cout << number << endl;
}
\end{lstlisting}

출력 결과:

\begin{verbatim}
1
3
5
7
9
\end{verbatim}

\subsection{\texttt{break}와 \texttt{continue} 비교}

\begin{center}
\begin{tabular}{>{\raggedright\arraybackslash}p{0.22\textwidth} >{\raggedright\arraybackslash}p{0.55\textwidth}}
\toprule
문법 & 동작 \\
\midrule
\texttt{break} & 반복문 전체를 즉시 종료한다. \\
\texttt{continue} & 현재 반복의 남은 부분만 건너뛰고 다음 반복을 수행한다. \\
\bottomrule
\end{tabular}
\end{center}

\subsection{예외를 먼저 처리하기}

\begin{lstlisting}
for (int count = 0; count < 5; count++)
{
    int number;
    cin >> number;

    if (number < 0)
    {
        cout << "Negative numbers are ignored." << endl;
        continue;
    }

    cout << "Square: " << number * number << endl;
}
\end{lstlisting}

음수가 입력되면 제곱 계산을 건너뛰고 다음 입력으로 넘어간다. 이처럼 예외적인 경우를 먼저 처리하면 정상 처리 코드의 중첩을 줄일 수 있다.

\begin{practicebox}
\textbf{파일: \texttt{07\_continue.cpp}}\\
1부터 30까지의 정수 중 3의 배수는 건너뛰고 나머지만 출력하라.
\end{practicebox}

% =================================================
\section{난수 생성의 기본 원리}
% =================================================

숫자 맞히기 게임을 만들려면 프로그램이 정답 숫자를 생성해야 한다. 전통적인 C++ 난수 생성에는 \texttt{rand()}, \texttt{srand()}, \texttt{time()}을 사용한다.

\begin{lstlisting}
#include <cstdlib>
#include <ctime>
\end{lstlisting}

\subsection{\texttt{rand()}}

\texttt{rand()}는 0 이상 \texttt{RAND\_MAX} 이하의 정수를 반환한다.

\begin{lstlisting}
cout << rand() << endl;
cout << rand() << endl;
cout << rand() << endl;
\end{lstlisting}

이 값들은 완전한 의미의 무작위가 아니라 정해진 계산 규칙에 따라 생성되는 의사 난수이다.

\subsection{seed와 \texttt{srand()}}

난수 생성기는 시작값을 기준으로 수열을 만든다. 이 시작값을 seed라고 하며, 같은 seed를 사용하면 같은 난수 수열이 생성된다.

\begin{lstlisting}
srand(10);
cout << rand() << endl;
\end{lstlisting}

프로그램을 실행할 때마다 seed를 바꾸기 위해 현재 시간을 사용할 수 있다.

\begin{lstlisting}
srand(time(0));
\end{lstlisting}

\texttt{time(0)}은 현재 시각에 따라 달라지는 값을 반환하므로 프로그램을 다시 실행할 때 다른 난수 수열이 생성될 가능성이 높다.

\begin{importantbox}
\texttt{srand(time(0))}은 일반적으로 \texttt{main()}의 시작 부분에서 한 번만 호출한다. 난수를 만들 때마다 반복 호출하면 같은 초 안에서 동일한 seed가 설정될 수 있다.
\end{importantbox}

\begin{practicebox}
\textbf{파일: \texttt{08\_random.cpp}}\\
난수 다섯 개를 출력하라. 먼저 \texttt{srand()} 없이 여러 번 실행하고, 이후 \texttt{srand(time(0))}을 추가하여 결과를 비교하라.
\end{practicebox}

% =================================================
\section{원하는 범위의 난수 만들기}
% =================================================

나머지 연산자 \texttt{\%}를 사용하면 난수의 범위를 제한할 수 있다.

\begin{lstlisting}
rand() % 100
\end{lstlisting}

100으로 나눈 나머지는 0부터 99까지이므로 위 표현식의 결과 범위는 다음과 같다.

\[
0 \le r \le 99
\]

1부터 100까지의 값을 원하면 1을 더한다.

\begin{lstlisting}
int number = rand() % 100 + 1;
\end{lstlisting}

\[
1 \le \texttt{number} \le 100
\]

\subsection{일반적인 범위 공식}

\begin{lstlisting}
rand() % (maximum - minimum + 1) + minimum
\end{lstlisting}

예를 들어 5부터 10까지의 정수를 생성하려면 다음과 같이 작성한다.

\begin{lstlisting}
int number = rand() % 6 + 5;
\end{lstlisting}

\begin{center}
\begin{tabular}{>{\raggedright\arraybackslash}p{0.45\textwidth} >{\raggedright\arraybackslash}p{0.28\textwidth}}
\toprule
표현식 & 생성 범위 \\
\midrule
\texttt{rand() \% 10} & 0--9 \\
\texttt{rand() \% 10 + 1} & 1--10 \\
\texttt{rand() \% 6 + 5} & 5--10 \\
\texttt{rand() \% 100 + 1} & 1--100 \\
\bottomrule
\end{tabular}
\end{center}

\begin{importantbox}
\texttt{rand() \% N} 방식은 학습용 프로그램에는 충분하지만, 엄밀하게 모든 값이 같은 확률로 생성된다고 보장되지는 않는다. 현대 C++에서는 \texttt{<random>} 라이브러리가 더 권장되지만, 여기서는 조건문과 반복문을 결합하는 데 집중한다.
\end{importantbox}

\begin{practicebox}
\textbf{파일: \texttt{09\_random\_range.cpp}}\\
0--9, 1--6, 10--20, $-5$--5 범위의 난수를 각각 생성하여 출력하라.
\end{practicebox}

% =================================================
\section{종합 실습: 숫자 맞히기 게임}
% =================================================

Day 2에서 학습한 조건문, \texttt{do--while}, 난수 생성을 결합하여 숫자 맞히기 게임을 작성한다.

\begin{enumerate}[leftmargin=1.8em]
    \item 프로그램이 1부터 100 사이의 정답을 생성한다.
    \item 사용자가 숫자를 입력한다.
    \item 정답보다 크면 \texttt{Too high!}를 출력한다.
    \item 정답보다 작으면 \texttt{Too low!}를 출력한다.
    \item 정답과 같으면 \texttt{Correct!}를 출력한다.
    \item 정답을 맞힐 때까지 반복한다.
\end{enumerate}

\subsection{완성 코드: \texttt{practice/p01\_number\_guessing.cpp}}

\begin{lstlisting}
#include <iostream>
#include <cstdlib>
#include <ctime>

using namespace std;

int main()
{
    srand(time(0));

    int answer = rand() % 100 + 1;
    int guess;
    int attempts = 0;

    cout << "==================================" << endl;
    cout << "       Number Guessing Game       " << endl;
    cout << "==================================" << endl;
    cout << "Guess a number from 1 to 100." << endl;

    do
    {
        cout << "Guess: ";
        cin >> guess;

        if (guess < 1 || guess > 100)
        {
            cout << "Enter a number from 1 to 100." << endl;
            continue;
        }

        attempts++;

        if (guess > answer)
        {
            cout << "Too high!" << endl;
        }
        else if (guess < answer)
        {
            cout << "Too low!" << endl;
        }
        else
        {
            cout << "Correct!" << endl;
        }
    }
    while (guess != answer);

    cout << "Attempts: " << attempts << endl;
    cout << "Game Over." << endl;

    return 0;
}
\end{lstlisting}

\subsection{실행 흐름 분석}

정답이 68이라고 가정하자. 사용자가 50을 입력하면 \texttt{guess < answer}가 참이므로 \texttt{Too low!}가 출력된다. 이후 \texttt{guess != answer}가 참이므로 반복한다.

75를 입력하면 \texttt{guess > answer}가 참이므로 \texttt{Too high!}가 출력된다. 마지막으로 68을 입력하면 앞의 두 조건이 모두 거짓이고 마지막 \texttt{else}가 실행된다. 이후 반복 조건이 거짓이 되어 게임이 종료된다.

입력값이 1부터 100 사이가 아니면 \texttt{continue}가 실행된다. 따라서 정답 비교와 시도 횟수 증가는 건너뛰고 다음 입력으로 이동한다.

% =================================================
\section{프로젝트 파일 구성}
% =================================================

\begin{verbatim}
cpp/
└── day02/
    ├── 01_if_else.cpp
    ├── 02_else_if.cpp
    ├── 03_nested_if.cpp
    ├── 04_switch.cpp
    ├── 05_do_while.cpp
    ├── 06_break.cpp
    ├── 07_continue.cpp
    ├── 08_random.cpp
    ├── 09_random_range.cpp
    ├── practice/
    │   └── p01_number_guessing.cpp
    ├── cpp_day02.tex
    ├── cpp_day02.pdf
    └── README.md
\end{verbatim}

개념별 예제 파일은 Day 2 루트에 번호를 붙여 분리하고, 종합 실습 파일은 \texttt{practice} 폴더에 \texttt{p01\_...}, \texttt{p02\_...} 형식으로 저장한다.

% =================================================
\section{핵심 정리}
% =================================================

\begin{importantbox}
\begin{itemize}[leftmargin=1.5em]
    \item \texttt{if--else}: 두 경우 중 하나를 선택한다.
    \item \texttt{else if}: 여러 조건을 위에서부터 순서대로 검사한다.
    \item \texttt{switch}: 하나의 값과 여러 상수를 비교한다.
    \item \texttt{do--while}: 본문을 먼저 실행하므로 최소 한 번 실행된다.
    \item \texttt{break}: 가장 가까운 반복문 또는 \texttt{switch}를 종료한다.
    \item \texttt{continue}: 현재 반복의 남은 코드를 건너뛴다.
    \item \texttt{srand(time(0))}: 현재 시간을 난수 seed로 설정한다.
    \item \texttt{rand() \% (max - min + 1) + min}: 지정 범위의 정수를 만든다.
\end{itemize}
\end{importantbox}

% =================================================
\section{체크 질문}
% =================================================

\begin{checkpointbox}
\begin{enumerate}[leftmargin=1.8em]
    \item \texttt{if--else}에서 두 블록이 동시에 실행될 수 있는가?
    \item \texttt{else if} 조건의 순서가 중요한 이유는 무엇인가?
    \item \texttt{switch}의 각 \texttt{case}에 \texttt{break}가 없으면 어떤 일이 발생하는가?
    \item \texttt{while}문과 \texttt{do--while}문의 가장 큰 차이는 무엇인가?
    \item \texttt{break}와 \texttt{continue}의 차이는 무엇인가?
    \item 중첩 반복문 안에서 \texttt{break}를 실행하면 몇 개의 반복문이 종료되는가?
    \item \texttt{rand() \% 10}의 결과 범위는 무엇인가?
    \item 1부터 10까지의 난수를 만들려면 어떤 표현식을 사용하는가?
    \item \texttt{srand(time(0))}은 왜 일반적으로 한 번만 호출하는가?
    \item 숫자 맞히기 게임에서 \texttt{do--while}문이 적합한 이유는 무엇인가?
\end{enumerate}
\end{checkpointbox}

% =================================================
\section{추가 연습 문제}
% =================================================

\begin{enumerate}[leftmargin=1.8em]
    \item 정수 하나를 입력받아 절댓값을 출력하라. 표준 라이브러리의 절댓값 함수는 사용하지 않는다.
    \item 서로 다른 정수 세 개를 입력받아 가장 큰 값을 출력하라.
    \item 정수 두 개와 메뉴 번호를 입력받는 간단한 계산기를 \texttt{switch}문으로 작성하라. 0으로 나누는 경우도 처리한다.
    \item 정답 비밀번호를 1234로 두고 올바른 값이 입력될 때까지 \texttt{do--while}문으로 반복하라.
    \item 1부터 50까지 출력하되 5의 배수는 \texttt{continue}로 건너뛰어라.
    \item 2 이상의 정수를 입력받고, 처음 발견한 약수를 출력한 뒤 \texttt{break}로 종료하라.
    \item 1부터 6 사이의 난수를 열 번 생성하여 주사위 굴리기를 구현하라.
    \item 20--30, $-10$--10, 100--200 범위의 난수 생성식을 각각 작성하라.
    \item 숫자 맞히기 게임에 최대 시도 횟수 7회를 추가하라.
    \item 게임이 끝난 뒤 \texttt{Play again? (y/n)}을 묻고, \texttt{y}이면 새 정답으로 다시 시작하도록 확장하라.
\end{enumerate}

\end{document}
